%------------------------------------------------------------------------------%
\documentclass[fleqn,utf8,aspectratio=169,12pt]{beamer}

\usepackage{calculo_varias_variaveis-1}

\title{Operadores Vetoriais}

%------------------------------------------------------------------------------%
\begin{document}

\addtitle

%------------------------------------------------------------------------------%
\begin{frame}{Operadores Vetoriais}
  As definições apresentadas nesta apresentação \emph{não} integram a ementa
  da disciplina. Sua inclusão tem caráter introdutório,
  com o objetivo de proporcionar um primeiro contato e promover a
  familiarização dos alunos com esses conceitos.
\end{frame}

%------------------------------------------------------------------------------%
\begin{frame}{Operadores Diferenciais Vetoriais}
  Seja $\mathbf{F} = (F_1, F_2, F_3)$ um campo vetorial (função de $\R^3$ em $\R^3$)

  \medskip
  $\phi$ um campo escalar (função de $\R^3$ em $\R$).

  \medskip
  suficientemente regulares
\end{frame}

%------------------------------------------------------------------------------%
\begin{frame}{Gradiente}
  \emph{Gradiente}
  \[
    \nabla \phi
    = \left( \D{\phi}{x},
             \D{\phi}{y},
             \D{\phi}{z}
      \right)
  \]

  Indica a direção e a taxa de maior variação de um campo escalar
\end{frame}

%------------------------------------------------------------------------------%
\begin{frame}{Divergente}
  \emph{Divergente}
  \[
    \nabla \cdot \mathbf{F}
      = \D{F_1}{x}
      + \D{F_2}{y}
      + \D{F_3}{z}
  \]

  Mede a intensidade com que um campo vetorial se expande ou se contrai
  em um ponto
\end{frame}

%------------------------------------------------------------------------------%
\begin{frame}{Rotacional}
\emph{Rotacional}
  \[
    \nabla \times \mathbf{F} =
    \left(
    \D{F_3}{y} - \D{F_2}{z},\quad
    \D{F_1}{z} - \D{F_3}{x},\quad
    \D{F_2}{x} - \D{F_1}{y}
    \right)
  \]

  Descreve a tendência de rotação ou circulação de um campo vetorial
  em torno de um ponto
\end{frame}

%------------------------------------------------------------------------------%
\begin{frame}{Laplaciano}
  \emph{Laplaciano}
  \[
    \nabla^2 \phi
    \;=\; \nabla \cdot (\nabla \phi)
    \;=\; \DS{\phi}{x} + \DS{\phi}{y} + \DS{\phi}{z}
  \]

  Expressa como um campo escalar difere, em um ponto,
  da média de seus valores nas proximidades
\end{frame}

%------------------------------------------------------------------------------%
\end{document}
%------------------------------------------------------------------------------%
